\documentclass[a4paper,12pt,twoside]{memoir}

% Castellano
\usepackage[spanish,es-tabla]{babel}
\selectlanguage{spanish}
\usepackage[utf8]{inputenc}
\usepackage[T1]{fontenc}
\usepackage{lmodern} % scalable font
\usepackage{microtype}
\usepackage{placeins}

\RequirePackage{booktabs}
\RequirePackage[table]{xcolor}
\RequirePackage{xtab}
\RequirePackage{multirow}

% Links
\PassOptionsToPackage{hyphens}{url}\usepackage[colorlinks]{hyperref}
\hypersetup{
	allcolors = {red}
}

% Ecuaciones
\usepackage{amsmath}

% Rutas de fichero / paquete
\newcommand{\ruta}[1]{{\sffamily #1}}

% Párrafos
\nonzeroparskip

% Huérfanas y viudas
\widowpenalty100000
\clubpenalty100000

% Evitar solapes en el header
\nouppercaseheads

% Imagenes
\usepackage{graphicx}
\newcommand{\imagen}[2]{
	\begin{figure}[!h]
		\centering
		\includegraphics[width=0.9\textwidth]{#1}
		\caption{#2}\label{fig:#1}
	\end{figure}
	\FloatBarrier
}

\newcommand{\imagenflotante}[2]{
	\begin{figure}%[!h]
		\centering
		\includegraphics[width=0.9\textwidth]{#1}
		\caption{#2}\label{fig:#1}
	\end{figure}
}



% El comando \figura nos permite insertar figuras comodamente, y utilizando
% siempre el mismo formato. Los parametros son:
% 1 -> Porcentaje del ancho de página que ocupará la figura (de 0 a 1)
% 2 --> Fichero de la imagen
% 3 --> Texto a pie de imagen
% 4 --> Etiqueta (label) para referencias
% 5 --> Opciones que queramos pasarle al \includegraphics
% 6 --> Opciones de posicionamiento a pasarle a \begin{figure}
\newcommand{\figuraConPosicion}[6]{%
  \setlength{\anchoFloat}{#1\textwidth}%
  \addtolength{\anchoFloat}{-4\fboxsep}%
  \setlength{\anchoFigura}{\anchoFloat}%
  \begin{figure}[#6]
    \begin{center}%
      \Ovalbox{%
        \begin{minipage}{\anchoFloat}%
          \begin{center}%
            \includegraphics[width=\anchoFigura,#5]{#2}%
            \caption{#3}%
            \label{#4}%
          \end{center}%
        \end{minipage}
      }%
    \end{center}%
  \end{figure}%
}

%
% Comando para incluir imágenes en formato apaisado (sin marco).
\newcommand{\figuraApaisadaSinMarco}[5]{%
  \begin{figure}%
    \begin{center}%
    \includegraphics[angle=90,height=#1\textheight,#5]{#2}%
    \caption{#3}%
    \label{#4}%
    \end{center}%
  \end{figure}%
}
% Para las tablas
\newcommand{\otoprule}{\midrule [\heavyrulewidth]}
%
% Nuevo comando para tablas pequeñas (menos de una página).
\newcommand{\tablaSmall}[5]{%
 \begin{table}
  \begin{center}
   \rowcolors {2}{gray!35}{}
   \begin{tabular}{#2}
    \toprule
    #4
    \otoprule
    #5
    \bottomrule
   \end{tabular}
   \caption{#1}
   \label{tabla:#3}
  \end{center}
 \end{table}
}

%
%Para el float H de tablaSmallSinColores
\usepackage{float}

%
% Nuevo comando para tablas pequeñas (menos de una página).
\newcommand{\tablaSmallSinColores}[5]{%
 \begin{table}[H]
  \begin{center}
   \begin{tabular}{#2}
    \toprule
    #4
    \otoprule
    #5
    \bottomrule
   \end{tabular}
   \caption{#1}
   \label{tabla:#3}
  \end{center}
 \end{table}
}

\newcommand{\tablaApaisadaSmall}[5]{%
\begin{landscape}
  \begin{table}
   \begin{center}
    \rowcolors {2}{gray!35}{}
    \begin{tabular}{#2}
     \toprule
     #4
     \otoprule
     #5
     \bottomrule
    \end{tabular}
    \caption{#1}
    \label{tabla:#3}
   \end{center}
  \end{table}
\end{landscape}
}

%
% Nuevo comando para tablas grandes con cabecera y filas alternas coloreadas en gris.
\newcommand{\tabla}[6]{%
  \begin{center}
    \tablefirsthead{
      \toprule
      #5
      \otoprule
    }
    \tablehead{
      \multicolumn{#3}{l}{\small\sl continúa desde la página anterior}\\
      \toprule
      #5
      \otoprule
    }
    \tabletail{
      \hline
      \multicolumn{#3}{r}{\small\sl continúa en la página siguiente}\\
    }
    \tablelasttail{
      \hline
    }
    \bottomcaption{#1}
    \rowcolors {2}{gray!35}{}
    \begin{xtabular}{#2}
      #6
      \bottomrule
    \end{xtabular}
    \label{tabla:#4}
  \end{center}
}

%
% Nuevo comando para tablas grandes con cabecera.
\newcommand{\tablaSinColores}[6]{%
  \begin{center}
    \tablefirsthead{
      \toprule
      #5
      \otoprule
    }
    \tablehead{
      \multicolumn{#3}{l}{\small\sl continúa desde la página anterior}\\
      \toprule
      #5
      \otoprule
    }
    \tabletail{
      \hline
      \multicolumn{#3}{r}{\small\sl continúa en la página siguiente}\\
    }
    \tablelasttail{
      \hline
    }
    \bottomcaption{#1}
    \begin{xtabular}{#2}
      #6
      \bottomrule
    \end{xtabular}
    \label{tabla:#4}
  \end{center}
}

%
% Nuevo comando para tablas grandes sin cabecera.
\newcommand{\tablaSinCabecera}[5]{%
  \begin{center}
    \tablefirsthead{
      \toprule
    }
    \tablehead{
      \multicolumn{#3}{l}{\small\sl continúa desde la página anterior}\\
      \hline
    }
    \tabletail{
      \hline
      \multicolumn{#3}{r}{\small\sl continúa en la página siguiente}\\
    }
    \tablelasttail{
      \hline
    }
    \bottomcaption{#1}
  \begin{xtabular}{#2}
    #5
   \bottomrule
  \end{xtabular}
  \label{tabla:#4}
  \end{center}
}



\definecolor{cgoLight}{HTML}{EEEEEE}
\definecolor{cgoExtralight}{HTML}{FFFFFF}

%
% Nuevo comando para tablas grandes sin cabecera.
\newcommand{\tablaSinCabeceraConBandas}[5]{%
  \begin{center}
    \tablefirsthead{
      \toprule
    }
    \tablehead{
      \multicolumn{#3}{l}{\small\sl continúa desde la página anterior}\\
      \hline
    }
    \tabletail{
      \hline
      \multicolumn{#3}{r}{\small\sl continúa en la página siguiente}\\
    }
    \tablelasttail{
      \hline
    }
    \bottomcaption{#1}
    \rowcolors[]{1}{cgoExtralight}{cgoLight}

  \begin{xtabular}{#2}
    #5
   \bottomrule
  \end{xtabular}
  \label{tabla:#4}
  \end{center}
}




\graphicspath{ {./img/} }

% Capítulos
\chapterstyle{bianchi}
\newcommand{\capitulo}[2]{
	\setcounter{chapter}{#1}
	\setcounter{section}{0}
	\setcounter{figure}{0}
	\setcounter{table}{0}
	\chapter*{#2}
	\addcontentsline{toc}{chapter}{#2}
	\markboth{#2}{#2}
}

% Apéndices
\renewcommand{\appendixname}{Apéndice}
\renewcommand*\cftappendixname{\appendixname}

\newcommand{\apendice}[1]{
	%\renewcommand{\thechapter}{A}
	\chapter{#1}
}

\renewcommand*\cftappendixname{\appendixname\ }

% Formato de portada
\makeatletter
\usepackage{xcolor}
\newcommand{\tutor}[1]{\def\@tutor{#1}}
\newcommand{\course}[1]{\def\@course{#1}}
\definecolor{cpardoBox}{HTML}{E6E6FF}
\def\maketitle{
  \null
  \thispagestyle{empty}
  % Cabecera ----------------
\noindent\includegraphics[width=\textwidth]{cabecera}\vspace{1cm}%
  \vfill
  % Título proyecto y escudo informática ----------------
  \colorbox{cpardoBox}{%
    \begin{minipage}{.8\textwidth}
      \vspace{.5cm}\Large
      \begin{center}
      \textbf{TFG del Grado en Ingeniería Informática}\vspace{.6cm}\\
      \textbf{\LARGE\@title{}}
      \end{center}
      \vspace{.2cm}
    \end{minipage}

  }%
  \hfill\begin{minipage}{.20\textwidth}
    \includegraphics[width=\textwidth]{escudoInfor}
  \end{minipage}
  \vfill
  % Datos de alumno, curso y tutores ------------------
  \begin{center}%
  {%
    \noindent\LARGE
    Presentado por \@author{}\\ 
    en Universidad de Burgos --- \@date{}\\
    Tutor: \@tutor{}\\
  }%
  \end{center}%
  \null
  \cleardoublepage
  }
\makeatother


% Datos de portada
\title{título del TFG \\Documentación Técnica}
\author{nombre alumno}
\tutor{nombre tutor}
\date{\today}

\begin{document}

\maketitle



\cleardoublepage



%%%%%%%%%%%%%%%%%%%%%%%%%%%%%%%%%%%%%%%%%%%%%%%%%%%%%%%%%%%%%%%%%%%%%%%%%%%%%%%%%%%%%%%%



\frontmatter


\clearpage

% Indices
\tableofcontents

\clearpage

\listoffigures

\clearpage

\listoftables

\clearpage

\mainmatter

\appendix

\apendice{Plan de Proyecto Software}

\section{Introducción}
En este apartado se va a analizar la planificación temporal y se va a estudiar la viabilidad económica y legal del proyecto.

\section{Planificación temporal}
Para la planficación temporal se han utilizado los Sprints de la metodología Scrum. Antes de empezar cada Sprint es necesario realizar una reunión para decidir la duración del mismo, el objetivo del mismo y las tareas que se van a llevar a cabo para cumplirlo. Una vez transcurrido ese tiempo se vuelve a celebrar una reunión y se actualizan los objetivos y tareas en función de lo realizado en el Sprint anterior. Y así sucesivamente hasta finalizar el desarrollo del proyecto software.

\subsection{Sprints}
Este apartado contiene los diferentes Sprints con las tareas que se han ido realizando en cada uno.

\subsubsection{Sprint Inicial}
\textbf{14/03/24 - 17/04/24}
\begin{itemize}
    \item Subir plantilla de la documentación a GitHub.
    \item Crear el fichero Python.
    \item Registro en Overleaf.
    \item Registro en Zube.
    \item Buscar información sobre Streamlit.
    \item Crear un prototipo de la aplicación.
    \item Añadir más funcionalidades al prototipo.
    \item Buscar información sobre el tratamiento de imágenes con OpenCV.
    \item Implementar las transformaciones a las imágenes.
\end{itemize}

\subsubsection{Sprint 2}
\textbf{18/04/24 - 10/05/24}
\begin{itemize}
    \item Añadir contenido inicial a la memoria.
    \item Cambiar canales RGB de las imágenes.
    \item Eliminar fondo de las imágenes.
    \item Determinar el porcentaje de la superficie de la hoja.
    \item Detectar el borde de las gotas.
    \item Determinar el porcentaje de la superficie de las gotas respecto a la hoja.
\end{itemize}

\subsubsection{Sprint 3}
\textbf{11/05/24 - 20/05/24}
\begin{itemize}
    \item Descargar imágenes hiperespectrales.
    \item Buscar información sobre las imágenes hiperespectrales y cómo trabajar con ellas.
    \item Crear método para cargar imágenes hiperespectrales.
    \item Crear método para visualizar las diferentes capas de la imagen hiperespectrales.
    \item Añadir documentación sobre las imágenes hiperespectrales.
\end{itemize}

\subsubsection{Sprint 4}
\textbf{21/05/24 - 28/05/24}
\begin{itemize}
    \item Crear pestaña de trinarización.
    \item Crear pestaña de visualización..
\end{itemize}

\subsubsection{Sprint 5}
\textbf{29/05/24 - 06/06/24}
\begin{itemize}
    \item Buscar librería de python para visualizar imágenes y poder editar parámetros.
    \item Poner la opción de elegir bandas en la pestaña Trinarizar.
    \item Mostrar porcentajes y guardarlos en un fichero .csv en la pestaña Trinarizar.
    \item Cambiar disposición de elementos en ambas pestañas.
    \item Cambiar el aspecto de los widgets de subir archivos y convertirlos en un botón.
    \item Desactivar la animación de carga de Hydralit.
    \item Crear el botón "Añadir Capas" en una ventana emergente.
\end{itemize}

\subsubsection{Sprint 6}
\textbf{07/06/24 - 28/06/24}
\begin{itemize}
    \item Añadir aspectos teóricos y herramientas utilizadas a la memoria.
\end{itemize}

\subsubsection{Sprint 7}
\textbf{29/06/24 - 03/07/24}
\begin{itemize}
    \item Implementar en la pestaña Visualizar el editor de imágenes con las capas.
    \item Carga de datos por defecto por fichero en la pestaña Trinarizar.
    \item Añadir información sobre el proyecto y su instalación a los 
    \item Crear la pestaña de Trinarizar por capas.
    anexos.
    \item Añadir transformaciones a la imagen trinarizada para eliminar ruido.
    \item Creación del fichero estilos.css.

\end{itemize}

\subsubsection{Sprint Final}
\textbf{04/07/24 - 09/07/24}
\begin{itemize}
    \item Terminar el resto apartados de la documentación.
    \item Desplegar la página web.
    \item Documentación interna del código.
\end{itemize}

\section{Estudio de viabilidad}
En este apartado se va realizar un análisis de la viabilidad del proyecto, teniendo en cuenta tanto el apartado económico, como el legal. 

\subsection{Viabilidad económica}
Antes de comenzar el desarrollo un proyecto software, es necesario hacer un cálculo de los gastos y de los posibles beneficios que nos vaya a aportar. Una vez calculados estos datos se podrá discernir si el proyecto es viable o no (desde un punto de vista puramente económico).
En los siguientes apartados se van a detallar los costes que supondrá el desarrollo de este proyecto. 

\subsubsection{Coste de Hardware}
El hardware utilizado para la realización de este proyecto ha sido un ordenador portátil y un ratón inalámbrico. 
\begin{itemize}
    \item Ordenador portátil MSI: 1200€
    \item Ratón inalámbrico Logitech: 50€
\end{itemize}
Se va a suponer una amortización a 5 años, por lo que el coste anual para el portátil sería de 1200€/5= 240€ y el del ratón sería de 50€/5=10€.

\subsubsection{Coste de Software}
El único software de pago que se ha utilizado en este proyecto es el sistema operativo del ordenador portátil, Windows 10 Pro.
\begin{itemize}
    \item Windows 10 Pro: 200€
\end{itemize}
Se va a suponer una amortización a 5 años, por lo que el coste anual para el software, Windows 10 Pro en este caso sería de 200€/5=40€.

\subsubsection{Coste de Personal}
Aquí vamos a considerar que estoy contratado por una empresa para la realización de este proyecto. Partiendo de esta base, se va a realizar una aproximación de los costes que les supondría contratar a un programador junior.
\begin{itemize}
    \item Sueldo anual de programador junior: 22000€
\end{itemize}
se va a suponer que se reparte el sueldo en 14 pagas y que se ha trabajado a media jornada, por lo que sería un sueldo mensual de 22000€/14 pagas = 1571,43€, pero como es a media jornada sería la mitad. Por lo que el sueldo mensual percibido corresponde a un total de 787,71€. El coste anual sería de 787,71€ X 12 meses = 9428,57€.

\subsubsection{Otros Costes}
En este apartado se incluyen gastos extras que han aparecido a lo largo de la realización del proyecto.
\begin{itemize}
    \item 3 Pendrives entrega: 30€
\end{itemize}
Estos gastos se amortizan en un año, por lo que valdrá 30€.

\subsubsection{Coste Total de un año de desarrollo}
Finalmente vamos a sumar todos los costes para ver cuánto dinero supondría la realización de este proyecto durante un año completo.

\begin{table}[h!]
	\centering
	\begin{tabular}{| l | r |}
		\toprule
		\textbf{Recurso} & \textbf{Coste Anual} \\ \midrule
		Windows 10 Pro & 40€\\
		Portátil &  240€ \\
            Ratón &  10€ \\
		Sueldo &  9428,57€\\
            Pendrives  &  30€\\
		\midrule
		\textbf{Total:} &  9748,57€ \\
		\bottomrule
	\end{tabular}
	\caption{Coste anual por recursos.}
	\label{CostResources}
\end{table}

Tras realizar los cálculos de gastos, no se ve viable el proyecto en términos económicos, ya que como no se consiga desarrollar un software muy potente y se pueda comercializar a los dueños de viñedos, para que reduzcan el uso de fertilizantes en sus campos, no se ve otra forma de conseguir ya no beneficio económico, sino simplemente cubrir los costes.

\subsection{Viabilidad legal}
En este apartado se van a revisar las licencias que poseen las diferentes librerías que se han utilizado en el proyecto, en qué consiste cada una, y también se va a explicar qué licencia encaja mejor con las características y finalidad de este proyecto.

\subsubsection{Licencias Software}
A continuación, se hace una breve descripción de las licencias que usan las librerías y se adjunta una tabla con las diferentes librerías, su versión y su licencia. 

\begin{itemize}
    \item \textbf{Licencia MIT}: La licencia MIT es una de las licencias de código abierto más permisivas. Permite el uso, modificación y distribución del software sin restricciones, siempre que se mantenga el aviso de derechos de autor original. No proporciona garantías ni responsabilidad a los autores por el uso del software.
	
    \item \textbf{Apache License 2.0}: La Apache License 2.0 es otra licencia de código abierto muy utilizada. Al igual que la licencia MIT, permite el uso, modificación y distribución del software de forma gratuita. Sin embargo, requiere mantener los avisos de derechos de autor y licencia tanto en el código original como en las versiones modificadas. También exime a los autores de cualquier responsabilidad o garantía sobre el software.
	
    \item \textbf{BSD-3-Clause License}: La BSD-3-Clause License es similar a las anteriores, permitiendo el uso, modificación y distribución del software de forma libre. A diferencia de la Licencia MIT, requiere que se indique si se han realizado cambios al trabajo original. Al igual que las otras, no ofrece garantías y exime a los autores de responsabilidad. Una característica destacada de esta licencia es su compatibilidad con otras licencias de código abierto, como la GPL.
	
\end{itemize}


\begin{table}[h!]
\centering
\begin{tabular}{|l|c|c|}
\hline \textbf{Librería} &\textbf{Versión} & \textbf{Licencia}\\ 
\hline hydralit & 1.0.14 &  Apache License 2.0\\
\hline numpy & 1.26.4 & Licencia BSD modificada\\
\hline opencv-python-headless & 4.10.0.84 & Apache License 2.0\\
\hline pandas & 2.2.1 & BSD 3-Clause License\\
\hline scikit-image & 0.24.0 & BSD 3-Clause License\\
\hline spectral & 0.23.1 & Licencia MIT\\
\hline streamlit & 1.36.0 & Apache License 2.0\\
\hline streamlit\_image\_zoom & 0.0.4 & Licencia MIT\\
\hline
\end{tabular}
\caption{Librerías utilizadas, su versión y la licencia que usan.}
\label{tab:librerias-licencias}
\end{table}

Tras analizar las diferentes licencias que utilizan las librerías y conocer sus características, se ha llegado a la conclusión de que la licencia MIT es la más adecuada para el proyecto. Esta elección se basa en que la aplicación se ha desarrollado para apoyar una investigación, y permitir que cualquier persona haga lo que desee con el código, siempre y cuando se cite al autor, facilita la colaboración y las mejoras del proyecto.


\apendice{Especificación de Requisitos}

\section{Introducción}

Una muestra de cómo podría ser una tabla de casos de uso:

% Caso de Uso 1 -> Consultar Experimentos.
\begin{table}[p]
	\centering
	\begin{tabularx}{\linewidth}{ p{0.21\columnwidth} p{0.71\columnwidth} }
		\toprule
		\textbf{CU-1}    & \textbf{Ejemplo de caso de uso}\\
		\toprule
		\textbf{Versión}              & 1.0    \\
		\textbf{Autor}                & Alumno \\
		\textbf{Requisitos asociados} & RF-xx, RF-xx \\
		\textbf{Descripción}          & La descripción del CU \\
		\textbf{Precondición}         & Precondiciones (podría haber más de una) \\
		\textbf{Acciones}             &
		\begin{enumerate}
			\def\labelenumi{\arabic{enumi}.}
			\tightlist
			\item Pasos del CU
			\item Pasos del CU (añadir tantos como sean necesarios)
		\end{enumerate}\\
		\textbf{Postcondición}        & Postcondiciones (podría haber más de una) \\
		\textbf{Excepciones}          & Excepciones \\
		\textbf{Importancia}          & Alta o Media o Baja... \\
		\bottomrule
	\end{tabularx}
	\caption{CU-1 Nombre del caso de uso.}
\end{table}

\section{Objetivos generales}

\section{Catálogo de requisitos}

\section{Especificación de requisitos}



\apendice{Especificación de diseño}

\section{Introducción}

\section{Diseño de datos}

\section{Diseño procedimental}

\section{Diseño arquitectónico}



\apendice{Documentación técnica de programación}

\section{Introducción}
En esta sección se explicará la estructura de directorios y archivos que componen el proyecto. También se enseñará a crear el entorno de desarrollo que necesita la aplicación en el manual del programador. Además, se detallará paso a paso cómo compilar, instalar y ejecutar la aplicación en diferentes sistemas operativos y utilizando diferentes métodos.

\section{Estructura de directorios}
El código fuente del proyecto se encuentra dentro de la carpeta \verb|/src|.

\subsection{src/}
En la carpeta src se encuentran una serie de ficheros:
\begin{itemize}
    \item \verb|bandas.xml|: Este archivo .xml contiene los valores predeterminados de bandas que toma el programa.
    \item \verb|estilos.css|: Este archivo .css contiene una serie de estilos utilizados en la aplicación.
    \item \verb|vitiscan.py|: En este archivo se encuentra el código fuente del proyecto.  
    \item \verb|requirements.txt|: Este archivo contiene las dependencias necesarias a instalar para poder compilar y ejecutar la aplicación.
    \item \verb|Dockerfile|: Este archivo es necesario para poder desplegar la aplicación en un contenedor Docker.
\end{itemize}

\subsection{src/.streamlit}
Esta carpeta contiene el fichero de configuración \verb|config.toml|

\section{Manual del programador}
En esta sección se explica cómo crear el entorno de desarrollo que necesita esta aplicación.

\subsection{Instalación de Python}

\subsubsection{GNU/Linux (Ubuntu)}
Para la instalación del entorno de desarrollo en GNU/Linux se ha decidido utilizar la distribución de Ubuntu, ya que es la distribución de Linux más popular.~\cite{ubuntu}

Lo primero que debemos hacer es instalar Python 3.11, para ello introduciremos  el siguiente comando en la consola:
\begin{verbatim}
    sudo apt install python3.11
\end{verbatim}
Para verificar la instalación se puede usar este comando:
\begin{verbatim}
    python3 --version
\end{verbatim}
Nos debería mostrar la versión de Python que hay instalada.

\subsubsection{Windows}
En el caso de Windows, debemos ir a la página oficial de Python y descargar el ejecutable de la versión 3.11.5 (\href{https://www.python.org/downloads/release/python-3115/}{enlace}).
Para verificar la instalación se puede usar este comando:
\begin{verbatim}
    python --version
\end{verbatim}
Nos debería mostrar la versión de Python que hay instalada.

\subsection{Creación y activación del entorno virtual}
Se recomienda utilizar un entorno virtual a la hora de crear un entorno de desarrollo en Python, ya que nos permite instalar las dependencias de forma independiente para cada proyecto.

\subsubsection{GNU/Linux (Ubuntu)}
Para crear el entorno virtual en Linux introduciremos el siguiente comando en la consola:
\begin{verbatim}
    python3 -m venv (nombre del entorno virtual)
\end{verbatim}

Posteriormente pasaremos a activar el entorno virtual introduciendo el siguiente comando en la consola:
\begin{verbatim}
    source (nombre del entorno virtual)/bin/activate
\end{verbatim}

Y con esto ya tendríamos todo listo para instalar las dependencias del proyecto.

\subsubsection{Windows}
Para crear el entorno virtual en Windows introduciremos el siguiente comando en la consola:
\begin{verbatim}
    python -m venv (nombre del entorno virtual)
\end{verbatim}

Posteriormente pasaremos a activar el entorno virtual introduciendo el siguiente comando en la consola:
\begin{verbatim}
    (nombre del entorno virtual)/Scripts/activate
\end{verbatim}

Y con esto ya tendríamos todo listo para instalar las dependencias del proyecto.

\subsection{Instalación de dependencias}
Esta sección es igual para los dos sistemas operativos por lo que procedemos a unificarlo. Para instalar las dependencias del proyecto hay que utilizar el siguiente comando en la consola:
\begin{verbatim}
    pip install -r requirements.txt
\end{verbatim}

Una vez hecho todo esto, ya tendríamos el entorno de desarrollo en funcionamiento. Ya se podría tanto modificar la aplicación y como ejecutar el proyecto en modo local, lo cuál se detalla en el siguiente apartado.

\label{sec:compilacion}
\section{Compilación, instalación y ejecución del proyecto}
El proyecto se puede ejecutar de varias maneras, que se procederán  a explicar a continuación. El código fuente se puede descargar del \href{https://github.com/dmp1002/TFG}{repositorio de GitHub} del proyecto. 

De todas formas se recomienda usar siguiente el comando, pero únicamente si se tiene el cliente de git instalado en el sistema:
\begin{verbatim}
    git clone https://github.com/dmp1002/TFG.git
\end{verbatim}
Este comando se deberá ejecutar en la ruta del entorno virtual creado previamente.

\subsection{Ejecución en local}
Suponiendo que se han realizado todos los pasos previamente explicados en el manual del programador, la ejecución en local no tiene ninguna dificultad.

Se debe ir a la ruta /src/ del proyecto y una vez allí ejecutar este comando en la consola:
\begin{verbatim}
    streamlit run vitiscan.py
\end{verbatim}

\subsection{Ejecución con Docker}
En este apartado se explica como desplegar la aplicación y cómo ejecutarla utilizando el software de creación de contenedores Docker.

\subsubsection{Instalación de Docker}
La instalación de Docker explicará únicamente para el sistema operativo Windows.

Lo primero de todo es instalar el Windows Subsystem for Linux, una funcionalidad de Windows que permite ejecutar distribuciones de Linux sin tener que usar una máquina virtual. Para instalarlo sólo hay que poner este comando en la consola Powershell.
\begin{verbatim}
    powershell wsl --install
\end{verbatim}

Una vez instalado te pedirá reiniciar. Al volver a encender el ordenador se abrirá una consola que te pedirá crear un usuario y contraseña para la version de Ubuntu que se ha instalado.

Antes de instalar Docker, debemos asegurarnos que en las características de Windows estén activadas dos opciones:
\begin{enumerate}
    \item El subsistema de Windows para Linux.
    \item La plataforma de máquina virtual.
\end{enumerate}

Una vez cumplidos estos requisitos, se pasa a descargar el instalador de Docker desde el siguiente enlace:

\href{https://docs.docker.com/desktop/install/windows-install/}{https://docs.docker.com/desktop/install/windows-install/}

Ya sólo falta ejecutarlo y una vez finalice la instalación reinciar el sistema, entonces ya estará listo para usarse.

\subsubsection{Construcción del contenedor Docker}
Para construir un contenedor Docker sólo necesitamos ir a la ruta donde tenemos instalado el proyecto y allí ejecutar este comando en la consola:
\begin{verbatim}
    docker build -t nombrecontenedor .
\end{verbatim}

Entonces se nos habrá creado el contenedor Docker donde podremos ejecutar el proyecto.
En nuestro caso lo llamaremos vitiscan para evitar confusiones.

\subsubsection{Ejecución del contenedor Docker}
Para ejecutar este contenedor se puede hacer por consola o directamente desde la interfaz del programa. Se recomienda hacerlo por comando, ya que es más sencillo. Sólo hay que ejecutar el siguiente comando en la misma ruta donde se haya creado el contenedor:
\begin{verbatim}
    docker run -p 8501:8501 nombrecontenedor
\end{verbatim}

Nosotros pondremos vitiscan para evitar confusiones.
Finalmente, se deberá abrir un navegador y acceder al siguiente enlace:
\href{http://localhost:8501/}{http://localhost:8501/}

En este proyecto se usa el puerto por defecto 8501, ya que así se ha configurado previamente en el fichero Dockerfile del proyecto.

\subsubsection{Exportación del contenedor Docker}
Con este comando se puede exportar el contenedor creado a una imagen Docker comprimida con extensión.tar con el siguiente comando:

\begin{verbatim}
    docker save -o nombrecontenedor.tar nombrecontenedor
\end{verbatim}


Así exportamos el contenedor actual a una imagen Docker comprimida .tar que se podrá cargar posteriormente en otro equipo y así poder ser ejecutada.
Nosotros llamaremos vitiscan tan a la imagen con al contenedor para evitar confusiones.

\section{Pruebas del sistema}
Aunque no se han llevado a cabo pruebas automatizadas para verificar el programa, se han realizado diferentes pruebas para asegurar que todas las funcionalidades implementadas se ejecutan correctamente.
\apendice{Documentación de usuario}

\section{Introducción}

\section{Requisitos de usuarios}

\section{Instalación}

\section{Manual del usuario}



\apendice{Anexo de sostenibilización curricular}

\section{Introducción}
Este anexo incluirá una reflexión personal del alumnado sobre los aspectos de la sostenibilidad que se abordan en el trabajo.
Se pueden incluir tantas subsecciones como sean necesarias con la intención de explicar las competencias de sostenibilidad adquiridas durante el alumnado y aplicadas al Trabajo de Fin de Grado.

Más información en el documento de la CRUE \url{https://www.crue.org/wp-content/uploads/2020/02/Directrices_Sosteniblidad_Crue2012.pdf}.

Este anexo tendrá una extensión comprendida entre 600 y 800 palabras.



\bibliographystyle{plain}
\bibliography{bibliografiaAnexos}

\end{document}
